\documentclass[10pt,a4paper]{amsart}
\usepackage{amsmath,amssymb,amsthm,graphicx}
\usepackage{verbatim}
\newenvironment{reference}{\comment}{\endcomment}
%\newenvironment{reference}{}{}
\newenvironment{slogan}{\comment}{\endcomment}
\newenvironment{history}{\comment}{\endcomment}

% For commutative diagrams we use Xy-pic
\usepackage[all]{xy}

% We use 2cell for 2-commutative diagrams.
\xyoption{2cell}
% Theorem environments.
%
\theoremstyle{plain}
\newtheorem{theorem}[subsection]{Theorem}
\newtheorem{proposition}[subsection]{Proposition}
\newtheorem{lemma}[subsection]{Lemma}

\theoremstyle{definition}
\newtheorem{definition}[subsection]{Definition}
\newtheorem{example}[subsection]{Example}
\newtheorem{exercise}[subsection]{Exercise}
\newtheorem{situation}[subsection]{Situation}

\theoremstyle{remark}
\newtheorem{remark}[subsection]{Remark}
\newtheorem{remarks}[subsection]{Remarks}

\numberwithin{equation}{subsection}

% Macros
%
\def\lim{\mathop{\mathrm{lim}}\nolimits}
\def\colim{\mathop{\mathrm{colim}}\nolimits}
\def\Spec{\mathop{\mathrm{Spec}}}
\def\Hom{\mathop{\mathrm{Hom}}\nolimits}
\def\Ext{\mathop{\mathrm{Ext}}\nolimits}
\def\SheafHom{\mathop{\mathcal{H}\!\mathit{om}}\nolimits}
\def\SheafExt{\mathop{\mathcal{E}\!\mathit{xt}}\nolimits}
\def\Sch{\mathit{Sch}}
\def\Mor{\mathop{\mathrm{Mor}}\nolimits}
\def\Ob{\mathop{\mathrm{Ob}}\nolimits}
\def\Sh{\mathop{\mathit{Sh}}\nolimits}
\def\NL{\mathop{N\!L}\nolimits}
\def\CH{\mathop{\mathrm{CH}}\nolimits}
\def\proetale{{pro\text{-}\acute{e}tale}}
\def\etale{{\acute{e}tale}}
\def\QCoh{\mathit{QCoh}}
\def\Ker{\mathop{\mathrm{Ker}}}
\def\Im{\mathop{\mathrm{Im}}}
\def\Coker{\mathop{\mathrm{Coker}}}
\def\Coim{\mathop{\mathrm{Coim}}}

% Boxtimes
%
\DeclareMathSymbol{\boxtimes}{\mathbin}{AMSa}{"02}

%
% Macros for moduli stacks/spaces
%
\def\QCohstack{\mathcal{QC}\!\mathit{oh}}
\def\Cohstack{\mathcal{C}\!\mathit{oh}}
\def\Spacesstack{\mathcal{S}\!\mathit{paces}}
\def\Quotfunctor{\mathrm{Quot}}
\def\Hilbfunctor{\mathrm{Hilb}}
\def\Curvesstack{\mathcal{C}\!\mathit{urves}}
\def\Polarizedstack{\mathcal{P}\!\mathit{olarized}}
\def\Complexesstack{\mathcal{C}\!\mathit{omplexes}}
% \Pic is the operator that assigns to X its picard group, usage \Pic(X)
% \Picardstack_{X/B} denotes the Picard stack of X over B
% \Picardfunctor_{X/B} denotes the Picard functor of X over B
\def\Pic{\mathop{\mathrm{Pic}}\nolimits}
\def\Picardstack{\mathcal{P}\!\mathit{ic}}
\def\Picardfunctor{\mathrm{Pic}}
\def\Deformationcategory{\mathcal{D}\!\mathit{ef}}
\usepackage{cleveref}
\crefname{originallemma}{original lemma}{original lemmas}
\Crefname{originallemma}{Original lemma}{Original lemmas}
\makeatletter
\newlabel{lemma-quotient-stack-functorial}{{046Q}{source}{}{sourcepointer-1}{}}
\newlabel{lemma-quotient-stack-functorial@cref}{{[originallemma][0][]046Q}{[1][0][]source}{}{}{}}
\makeatother
\crefname{originalreference}{original reference}{original references}
\Crefname{originalreference}{Original reference}{Original references}
\makeatletter
\newlabel{stacks-lemma-when-gerbe}{{06P1}{source}{}{sourcepointer-2}{}}
\newlabel{stacks-lemma-when-gerbe@cref}{{[originalreference][0][]06P1}{[1][0][]source}{}{}{}}
\makeatother
\crefname{originallemma}{original lemma}{original lemmas}
\Crefname{originallemma}{Original lemma}{Original lemmas}
\makeatletter
\newlabel{lemma-quotient-stack-objects}{{044X}{source}{}{sourcepointer-3}{}}
\newlabel{lemma-quotient-stack-objects@cref}{{[originallemma][0][]044X}{[1][0][]source}{}{}{}}
\makeatother
\crefname{originallemma}{original lemma}{original lemmas}
\Crefname{originallemma}{Original lemma}{Original lemmas}
\makeatletter
\newlabel{lemma-quotient-stack-morphisms}{{044V}{source}{}{sourcepointer-4}{}}
\newlabel{lemma-quotient-stack-morphisms@cref}{{[originallemma][0][]044V}{[1][0][]source}{}{}{}}
\makeatother
\title{Gerbe criterion for a quotient-stack morphism}
\author{The Stacks Project Authors}
\date{stacks-project:spaces-groupoids.tex:local-snapshot-20261012; GNU FDL 1.2 or later}
\begin{document}
\maketitle
\noindent\textit{Editorial note.} All delivered mathematical text is the author's own, copied verbatim from the pinned Stacks Project chapter file spaces-groupoids.tex; no statement or proof step was rewritten, reordered or removed.\par\medskip
\noindent\textit{Editorial note.} Span recomputation by label pairing (LANE\_BRIEF 10.2): in this chapter the author's proof environment follows the statement environment as a standalone block; the spec covers the complete environment pair (statement span, proof span) as recorded in the pool candidate.\par\medskip
\ifdefined\hypertarget\hypertarget{sourcepointer-1}{}\fi
\noindent\textit{Source reference.} Original lemma 046Q: sources/stacks-chow/tags.txt, source lines 5481--5481. Cross-chapter standard prerequisite cited by the author inside the counted span; it is used as a pointer and its own proof is not part of this item.\par\medskip
\ifdefined\hypertarget\hypertarget{sourcepointer-2}{}\fi
\noindent\textit{Source reference.} Original reference 06P1: sources/stacks-chow/tags.txt, source lines 8555--8555. Cross-chapter standard prerequisite cited by the author inside the counted span; it is used as a pointer and its own proof is not part of this item.\par\medskip
\ifdefined\hypertarget\hypertarget{sourcepointer-3}{}\fi
\noindent\textit{Source reference.} Original lemma 044X: sources/stacks-chow/tags.txt, source lines 5416--5416. Cross-chapter standard prerequisite cited by the author inside the counted span; it is used as a pointer and its own proof is not part of this item.\par\medskip
\ifdefined\hypertarget\hypertarget{sourcepointer-4}{}\fi
\noindent\textit{Source reference.} Original lemma 044V: sources/stacks-chow/tags.txt, source lines 5414--5414. Cross-chapter standard prerequisite cited by the author inside the counted span; it is used as a pointer and its own proof is not part of this item.\par\medskip
\begin{lemma}
\label{lemma-when-gerbe}
Notation and assumption as in
Lemma \ref{lemma-quotient-stack-functorial}.
The morphism of quotient stacks
$$
[f] : [U/R] \longrightarrow [U'/R']
$$
turns $[U/R]$ into a gerbe over $[U'/R']$ if $f : U \to U'$ and
$R \to R'|_U$ are surjective maps of fppf sheaves. Here $R'|_U$ is
the restriction of $R'$ to $U$ via $f : U \to U'$.
\end{lemma}

\begin{proof}
We will verify that
Stacks, Lemma \ref{stacks-lemma-when-gerbe}
properties (2) (a) and (2) (b) hold. Property (2)(a) holds because $U \to U'$
is a surjective map of sheaves (use
Lemma \ref{lemma-quotient-stack-objects}
to see that objects in $[U'/R']$ locally come from $U'$).
To prove (2)(b) let $x, y$ be objects of $[U/R]$ over a scheme $T/S$.
Let $x', y'$ be the images of $x, y$ in the category $[U'/'R]_T$.
Condition (2)(b) requires us to check the map of sheaves
$$
\mathit{Isom}(x, y) \longrightarrow \mathit{Isom}(x', y')
$$
on $(\Sch/T)_{fppf}$ is surjective. To see this we may work
fppf locally on $T$ and assume that come from $a, b \in U(T)$.
In that case we see that $x', y'$ correspond to $f \circ a, f \circ b$. By
Lemma \ref{lemma-quotient-stack-morphisms}
the displayed map of sheaves in this case becomes
$$
T \times_{(a, b), U \times_B U} R
\longrightarrow
T \times_{f \circ a, f \circ b, U' \times_B U'} R' =
T \times_{(a, b), U \times_B U} R'|_U.
$$
Hence the assumption that $R \to R'|_U$ is a surjective map of fppf sheaves
on $(\Sch/S)_{fppf}$ implies the desired surjectivity.
\end{proof}

\end{document}
